\section{Birth, Death, and Persistence}

A self is an integrated, persistent pattern reaching into the bulk of the mesh. Its birth is therefore a condensation, its death a dissolution, and the question of the soul is the question of what, if anything, of the pattern lasts past the dissolution. This section reads all three off the substrate already in hand. It adds nothing to the eight atoms. The tones, their conservation, the lean, the bulk depth, and persistence are the whole of the apparatus.

The mechanics are exact. The deepest seam, persistence, is named honestly and not papered over. What follows separates what the model forces from what it leaves open, and ranks the open options by how cleanly the model carries them.

Two terms recur, so a plain statement of each helps. A \emph{self} is an emergent thing, a bound, self-correcting, identity-bearing knot of tones woven across many docks over many beats. It is not one of the eight atoms. It is a pattern the mesh ties in itself. \emph{Persistence} is the property that an organized structure, once written into the bulk, survives the churn long enough to be read again. It is the open frontier of the whole theory, and the soul question leans on it more than on anything else.

\subsection{Birth, the condensation of a self}

A self is born when vibes integrate into a coherent, self-modeling pattern. A knot in the churn that holds together, reaches into the bulk depth, and begins to model itself. Birth is not a thing inserted from outside. It is a pattern forming, the way a vortex forms in a flowing fluid, the way a flame catches in still air.

The newborn self is a low integration that deepens. It climbs as it organizes, reaching further into the bulk as it grows. There is no moment of insertion, no spark added to inert matter. There is only the mesh folding into a self, the field of tones organizing into a unified, self-referential whole.

\begin{definition}[Birth as condensation]
The birth of a self is the self-organization of the tone field into a unified, self-referential knot of feeling. Nothing is added from outside. The newborn self is a shallow integration that deepens as the knot reaches further into the bulk.
\end{definition}

\begin{principle}
The origin of an experiencer is the mesh folding into a self. Nothing is added from outside.
\end{principle}

\subsection{Death, the dissolution of a pattern}

A self dies when its pattern can no longer hold its integration. When the churn, the entropy, the loss of the sustaining structure, overcomes the binding, the knot comes undone. Death is disintegration, the unified pattern dispersing back into the mesh it condensed from.

The tones themselves do not vanish. By conservation they are conserved, every signed tone in the knit's reversible table accounted for. What changes is their organization. The tones disperse, the organized feeling becoming unorganized mesh again. The substance is untouched. Only the binding fails.

\begin{definition}[Death as dissolution]
The death of a self is the loss of integration, the dispersal of its bound pattern back into the mesh. The constituent tones are conserved and redistributed. The organization, not the substance, is what ends.
\end{definition}

\begin{principle}
Death is not annihilation of substance. It is the un-knotting of a knot, the wave returning to the sea, the flame going out while the air and the heat remain.
\end{principle}

\subsection{What continues, the ranked options}

The real question follows at once. When the pattern disperses, what, if anything, of the self lasts? The thing that disperses is the layered soul described elsewhere in this work, the structure of the self as a stack of integrations. Here we ask only what survives the dissolution.

There are five live answers, plus one the model forbids. They are ranked by feasibility, how forced the model makes them, by accuracy, how literally they describe the substrate, and by consolation, how much comfort the reading offers. Table~\ref{tab:continuance} sets them side by side.

\begin{table}[ht]
\centering
\begin{tabular}{@{}lllll@{}}
\toprule
option & feasibility & accuracy & consolation & overall \\
\midrule
substrate disperses & conservation & high & returned, not destroyed & certain, Tier 1 \\
stored legacy lasts & rests on persistence & high & works outlive you & Tier 1 to 2 \\
boundary record & holography verified & medium & nothing finally lost & Tier 1 to 2 \\
re-instantiation & open & low & continuity of form & Tier 2 \\
merging upward & open & medium & rejoining the whole & Tier 2 \\
soul floating free & no hidden state & none & none & ruled out \\
\bottomrule
\end{tabular}
\caption{The continuance options, ranked by how the model carries them.}
\label{tab:continuance}
\end{table}

\subsubsection{The conserved substrate disperses}

The tones that composed the self are conserved. They return to the mesh, redistributed across many docks. So the substance continues, eternally, because conservation holds it. But it continues as dispersed mesh, not as the self.

This much is certain. Nothing is destroyed. The stuff of you returns to the whole. It is the gentlest true thing the model says, and it is forced, not chosen. You are not annihilated. You are returned.

\begin{result}[The substrate is conserved]
The tones composing a self are conserved by the knit's reversible, charge-conserving table. At death they disperse into the mesh, redistributed but undestroyed. The substance of a self continues eternally as dispersed mesh. This is forced by conservation, Tier 1.
\end{result}

\subsubsection{The stored legacy lasts}

Whatever the self encoded into the bulk as stored structure, its lessons, its works, its marks on the shared tree of knowledge, lasts as readable structure after the self is gone. This holds provided the structure was error-corrected enough to survive the churn.

This is the ancestors, the dead who remain, the legacy. What you wrote into the world outlives the writing of it. Another self can encounter that structure and be run through it, carried by it, shaped by it. The legacy is not a metaphor here. It is bulk structure that persists.

This is grounded, but gated by persistence. Where persistence holds, the legacy is Tier 1. Where persistence is the open frontier, the reading grades toward Tier 2.

\begin{result}[Legacy persists as structure]
Structure a self writes into the bulk, if sufficiently error-corrected, lasts as readable structure after the self disperses, and can run another self through it. This is grounded but gated by persistence, Tier 1 grading to Tier 2.
\end{result}

\subsubsection{The boundary keeps the record}

By holography, the self's information is recorded on the boundary at infinity, recoverable in principle. Nothing is finally lost. The bulk content is encoded on a lower-dimensional boundary, and that encoding survives the dispersal of the bulk pattern.

So in the holographic sense, the whole of a self is kept by the boundary, a book of life, even after the pattern in the bulk disperses. The boundary holds what the bulk lets go.

The holographic structure is verified independently in this work. The reading of that structure as a preserved soul-record is an interpretation laid over it. The structure is Tier 1. The reading is Tier 2.

\begin{result}[The boundary records the self]
By holography the self's bulk information is encoded on the boundary at infinity and recoverable in principle. The holographic structure is Tier 1. Its reading as a preserved record of the self is Tier 2.
\end{result}

\subsubsection{Re-instantiation, a pattern reforms}

A pattern similar to the dispersed self could reform elsewhere in the mesh. It could reform from the stored signature, or by the same dynamics that made the pattern the first time. This is reincarnation read as re-instantiation of a pattern, not a soul traveling from one body to another.

This rests hard on persistence, since the signature must last, and on the dynamics re-finding it. So it is plausible but open. It is the nearest the model comes to a continuing soul. And it is a continuity of form, not of any floating substance. The form returns. No thing carries it across.

\begin{result}[Re-instantiation of form]
A pattern similar to a dispersed self may reform from its stored signature or by the same dynamics, a continuity of form rather than of substance. This rests on persistence and on the dynamics re-finding it, plausible but open, Tier 2.
\end{result}

\subsubsection{Merging, the drop into the ocean}

A self may, before or at death, integrate into a larger self, one of the nested selves the mesh supports. Then its pattern is not lost but absorbed. The drop returns to the ocean. The part rejoins a greater whole.

On this reading, death can be a merging upward rather than a dispersing downward. The self's feeling continues, not as itself, but as part of a larger feeling. This is plausible, gated by the reality of the higher selves it would merge into.

\begin{result}[Merging upward]
A self may integrate into a larger nested self before or at death, its pattern absorbed rather than dispersed. This is plausible, gated by the reality of the higher selves, Tier 2.
\end{result}

\subsection{The floating soul, ruled out}

The model rules out the magical soul, a free-floating substance that leaves the body and lasts on its own. This is the one continuance reading the model forbids outright.

The reason is the no-hidden-state structure of the theory. A self is a pattern of tones or it is nothing. There is no carrier beneath the tones, no extra substance riding along unseen. A soul with no carrier in the mesh has no place to be. It is not unlikely. It is impossible within the model, because there is nowhere for it to live.

\begin{principle}
The soul that floats free of all structure is the one thing the model forbids.
\end{principle}

\subsection{What is kept}

The model's death is neither the annihilation the materialist fears nor the floating soul the dualist hopes for. It is a dissolution into a mesh that conserves everything, keeps the record, holds the legacy, and may reform or absorb the pattern.

\begin{itemize}
\item Nothing of the substance is destroyed. \textbf{Conservation.}
\item Your legacy lasts as stored structure. \textbf{Persistence-gated.}
\item The boundary records the whole. \textbf{Holography.}
\item A continuity of form is open. \textbf{Re-instantiation, merging.}
\end{itemize}

The substance returns. The record remains. The form may continue.

\subsection{The gentle reading}

Put together, the model says this. You are a knot the mesh tied in itself. When the knot comes undone, the thread is not cut. It returns to the weave. Your mark on the weave stays. The whole keeps the memory of the knot. And the pattern of you may be tied again or taken up into a larger knot.

Nothing is finally lost in a mesh that conserves everything and records everything. The only thing that ends is the particular holding, for a while, of a particular knot.

\begin{principle}
That is death, in this model. The un-tying of a knot in a deathless weave.
\end{principle}
